\section{관련 연구}

\subsection{PQC TLS 전환 연구}

Open Quantum Safe(OQS) 프로젝트는 양자 내성 암호 전환을 지원하기 위한 오픈소스 프로젝트로, PQC 알고리즘 구현체인 liboqs와 OpenSSL~3에서 PQC 알고리즘을 사용할 수 있도록 하는 oqs-provider를 제공한다~[인용 추가].
OQS는 OpenSSL~3 provider 기반으로 TLS~1.3, X.509, S/MIME에서 PQC 키 교환과 인증을 실험할 수 있는 환경을 제공하며, 이를 활용하여 Nginx와 curl 기반 TLS 실험 구성을 만들 수 있다~[인용 추가].
TLS~1.3에서 Hybrid 키 교환을 지원하기 위한 X25519MLKEM768 등의 그룹은 IETF Internet-Draft로 제안되어 있으며~[인용 추가], 대형 웹 서비스 사업자들은 실제 트래픽 환경에서 hybrid post-quantum TLS의 배포 가능성을 검증하는 실험을 수행해 왔다~[인용 추가].
NIST는 2024년 ML-KEM(FIPS~203), ML-DSA(FIPS~204), SLH-DSA(FIPS~205)를 PQC 표준으로 확정하였다~[인용 추가].

그러나 기존 연구는 PQC TLS 실험 기반을 제공하는 데 초점을 두며, 단계적 전환 정책의 코드화 및 CI/CD 파이프라인을 통한 자동 배포·검증 구조는 별도로 설계해야 한다.
본 연구는 OQS 스택을 기반으로 하되, Classical ECC에서 Hybrid PQC를 거쳐 Advanced Hybrid PQC로의 단계적 전환을 파이프라인이 자동으로 집행하는 운영 흐름을 설계·구현한다는 점에서 차별화된다.

\subsection{DevSecOps 및 CI/CD 보안 통합}

DevSecOps는 개발·보안·운영을 통합하여 보안을 CI/CD 파이프라인의 일부로 내재화하는 방법론이다~[인용 추가].
Trivy, Gitleaks, Semgrep과 같은 도구는 CI/CD 파이프라인에서 취약점 스캔, 시크릿 탐지, 정적 분석을 자동화하는 데 활용된다~[인용 추가].
또한 Open Policy Agent(OPA)로 대표되는 Policy as Code(PaC) 접근은 인프라 및 보안 정책을 코드로 명세하고 자동으로 집행하는 구조를 제공한다~[인용 추가].

그러나 기존 DevSecOps 접근은 일반적인 취약점 탐지와 배포 차단에 집중하며, TLS 암호 정책을 단계별로 정의하고 PQC 전환 진척도를 파이프라인이 집행하는 구조는 상대적으로 충분히 다루지 않는다.
본 연구는 TLS 설정의 암호 알고리즘 준수 여부를 Policy as Code로 강제하고, 정책 위반 시 배포를 자동 차단하는 Security Gate를 파이프라인에 통합한다는 점에서 차별화된다.

\subsection{CBOM 및 암호화 자산 관리}

CBOM(Cryptographic Bill of Materials)은 시스템 내 암호 알고리즘, 키, 인증서 등 암호 자산과 그 의존 관계를 기술하기 위한 객체 모델이며, CycloneDX는 CBOM 표현을 지원한다~[인용 추가].
CBOM을 활용하면 시스템에서 사용 중인 암호 알고리즘과 인증서 구성을 식별하고, 양자 취약 암호 자산의 목록을 구성할 수 있다.
PQC 전환 지침들은 일반적으로 전환 대상 암호 자산의 식별과 목록화를 선행 단계로 제시한다~[인용 추가].

그러나 기존 접근은 CBOM 생성 자체에 집중하며, 전환 전후의 CBOM을 자동으로 비교하여 마이그레이션 진척도를 BASELINE·IMPROVED·UNCHANGED·REGRESSED로 판정하는 구조는 충분히 다루지 않는다.
본 연구는 CycloneDX~1.7 기반 CBOM 생성과 실행 이력 간 자동 Diff를 구현하여, 전환 진척도를 객관적으로 추적하고 회귀를 즉시 감지하는 구조를 제시한다.

\subsection{LLM 기반 보안 코드 마이그레이션}

대형 언어 모델(LLM)을 활용한 코드 자동 변환 연구는 일반 코드 마이그레이션, 취약 코드 자동 수정 등 다양한 분야에서 진행되고 있다~[인용 추가].
또한 LLM 기반 보안 코드 생성 및 분석 연구들은 LLM이 암호화 관련 API를 잘못 사용하거나 보안 취약점을 포함한 코드를 생성할 수 있음을 보고하였다~[인용 추가].
한편 정적 분석과 LLM을 결합하여 코드 수정을 자동화하려는 시도도 존재한다~[인용 추가].

그러나 기존 연구는 주로 일반 코드 생성, 취약점 수정, 암호 API 오용 탐지에 초점을 두며, RSA에서 ML-KEM/ML-DSA로의 PQC 전환 코드 마이그레이션을 대상으로 정적 검증과 의미적 실행 검증을 분리 측정한 연구는 제한적이다.
본 연구는 PQC 전환 전용 합성 패턴 N=16을 설계하고, 4단계 검증 계층(L1--L4)을 적용한 80 trial 정량 평가를 통해 정적 AST 검증과 의미적 실행 검증의 역할을 분리 측정한다.